\chapter{相关理论与关键概念}

\section{属性基加密的基本思想}

属性基加密（ABE）的核心思想是把“解密权利”从具体身份扩展为属性集合和访问结构之间的匹配关系。传统公钥加密要求加密者明确知道接收者的公钥；ABE 则允许系统使用“属性”和“访问策略”描述授权条件。需要区分的是，在 ABE 的一般定义中，访问策略不一定都放在密文中：CP-ABE 将访问策略嵌入密文、用户私钥携带属性集合；KP-ABE 则通常由密文携带属性集合、用户私钥携带访问策略。因此，更准确地说，ABE 的解密能力来自密文侧信息与私钥侧信息是否满足同一访问结构，而不是单纯由某一方独立决定。

从访问控制角度看，ABE 将密码学和权限管理结合起来。一般的访问控制系统依赖服务器或数据库在访问时进行权限判断，而 ABE 把访问结构嵌入密钥生成和密文生成过程，使授权条件成为解密算法能否成功的一部分。即使密文被存储在不完全可信的云端，服务器也不能简单绕过访问条件，因为恢复明文需要私钥分量、密文分量和访问结构之间满足相应的代数关系。

落实到本文采用的 CP-ABE 场景，密文携带访问策略，用户私钥携带属性集合。Bethencourt、Sahai 和 Waters 提出的 CP-ABE 方案~[2]给出的直接动机正是“不完全可信存储环境中的加密访问控制”。在该模型中，数据拥有者可以在加密时指定访问策略，而不必预先知道所有接收者的身份；当且仅当用户属性集合满足密文中的策略时，用户才能通过正常解密流程恢复明文。这一点与本项目的使用方式一致：加密端输入访问控制树，解密端依赖用户私钥中的属性集合。

这种思想与传统访问控制相比有一个重要变化：访问判断不再只是软件系统中的一个 if 条件，而是变成密码算法本身的一部分。对于本文实现的 CP-ABE 而言，如果用户属性集合不能满足密文策略，即使用户得到密文文件和程序，也无法通过正常解密流程恢复明文。因此，在本实验中，策略表达、密钥分量和解密方程必须保持一致，否则就会出现“策略看起来正确但实际无法判定”的问题。

\section{CP-ABE 与 KP-ABE}

ABE 通常分为密文策略 ABE（CP-ABE）和密钥策略 ABE（KP-ABE）。二者的主要区别在于策略放在哪里。

\begin{table}[H]
\centering
\caption{CP-ABE 与 KP-ABE 的对比}
\begin{tabular}{p{0.22\textwidth} p{0.32\textwidth} p{0.32\textwidth}}
\toprule
方案类型 & 密文携带 & 私钥携带 \\
\midrule
CP-ABE & 访问策略，例如 $(A \wedge B)\vee C$ & 用户属性集合，例如 $\{A,B\}$ \\
KP-ABE & 属性集合 & 访问策略 \\
\bottomrule
\end{tabular}
\end{table}

本实验实现的是 CP-ABE。其使用方式更符合数据拥有者主导授权的场景：加密者在加密时写出访问策略，授权中心为用户按照实际身份或角色发放属性私钥。用户能否解密由“用户属性集合是否满足密文策略”决定。

选择 CP-ABE 还有一个工程上的好处：策略随密文保存，密文 JSON 中保留策略树、矩阵和行属性映射。这样报告中的理论推导、代码中的矩阵构造和测试中的解密结果可以对应到同一份密文文件，便于检查和复现实验过程。

实验三课件~[1]对 Waters 2011 方案~[3]的描述也采用了这一结构：加密算法输入公共参数、消息和访问结构 $(M,\rho)$，私钥生成算法根据用户属性集合生成密钥，解密算法再根据满足矩阵的行集合求重构常数。本项目保留这个 CP-ABE 方向，而没有实现 KP-ABE，是因为期末题目明确要求“属性基加密方案”并强调“使用访问控制树构建 LSSS 矩阵”，与 Waters 2011 式密文策略流程更直接对应。

\section{访问控制树}

访问控制树是一种直观的访问策略表示方法。树的叶子节点是属性字符串，例如 A、B、C、\verb|role:admin|；内部节点是阈值门 $k$-of-$n$，表示 $n$ 个子节点中至少有 $k$ 个满足。AND 和 OR 是阈值门的特例：

\begin{itemize}
    \item AND 门：$n$-of-$n$，所有子节点都必须满足。
    \item OR 门：$1$-of-$n$，至少一个子节点满足。
    \item 一般阈值门：$k$-of-$n$，至少 $k$ 个子节点满足。
\end{itemize}

本项目使用的示例策略为：

\[
    (A \wedge B)\vee C.
\]

它表示用户如果同时拥有属性 A 和 B，则可以解密；如果单独拥有属性 C，也可以解密；如果只有属性 A 或只有属性 B，则不能解密。

从访问控制语义看，本项目处理的是单调访问结构。所谓单调，是指用户在已经满足策略的属性集合中继续增加属性，不会使其从“满足”变为“不满足”。AND、OR 和阈值门都属于单调结构，而 NOT 条件不属于本实验实现范围。选择单调访问树与 CP-ABE 中常见的 LSSS 访问结构相一致，也便于后续把树上份额传播过程转换为矩阵行。

使用阈值门统一表示 AND 和 OR，可以减少实现中的特殊分支。系统只需要处理一般的 $k$-of-$n$ 节点，AND 被看作 $n$-of-$n$，OR 被看作 $1$-of-$n$。这种统一表达也方便后续扩展，例如测试中加入的 $2$-of-$3(D,E,F)$ 策略并不需要重新设计解密逻辑。换言之，访问树中的内部节点不需要被预先限制为二叉布尔门，只要给定阈值 $k$ 和子节点列表，就可以按照统一规则参与 LSSS 矩阵构造。

在程序表示上，叶子节点只保存属性名，内部节点保存 \verb|k| 和 \verb|children|。例如单个属性可以写成 \verb|{"attr":"A"}|，二选一或二者都满足的条件则可以写成带有子节点数组的阈值门对象。这样的 JSON 结构接近矩阵构造所需的递归输入：遇到叶子节点时生成矩阵行，遇到内部节点时继续向子节点分发份额。因此，策略文件既是用户可检查的访问规则，也是程序生成 LSSS 矩阵的直接输入。

访问树还允许同一个属性在不同叶子中出现多次。此时 LSSS 矩阵中可能有多行映射到同一个属性，即 $\rho$ 不一定是单射。Waters~[3]的 LSSS 表述允许这种情况，本项目的重复属性测试也验证了解密阶段不能简单地按属性名去重，而应保留所有满足 $\rho(i)\in S$ 的矩阵行。这个细节对访问树的一般性很重要，因为真实策略中同一属性可能出现在多个子策略位置。

Bethencourt 等人~[2]使用单调访问树描述密文策略，其中内部节点可以看作阈值门，叶子节点表示属性；Liu、Cao 和 Wong~[4]进一步从工程角度强调，AND 和 OR 都是一般阈值门的特例，直接支持阈值访问树可以避免把 $k$-of-$n$ 策略先展开成很大的布尔公式。本项目选择在 JSON 策略中直接保存 \verb|k| 和 \verb|children|，正是为了让 AND、OR 和一般阈值门共享同一种表示。

\section{线性秘密共享方案}

线性秘密共享方案（LSSS）可以用矩阵 $M$ 和行标号映射 $\rho$ 表示。设矩阵

\[
    M \in \mathbb{Z}_p^{\ell \times n},
\]

其中 $\ell$ 是矩阵行数，通常对应访问策略中的叶子节点数量；$n$ 是列数，对应秘密和若干随机变量的线性组合维度。映射

\[
    \rho:\{1,\dots,\ell\}\rightarrow \mathcal{U}
\]

将每一行映射到一个属性。

分享阶段选择向量

\[
    \mathbf{v}=(s,r_2,\dots,r_n),
\]

其中 $s$ 是需要共享的秘密，其余分量是随机数。第 $i$ 行得到的份额为：

\[
    \lambda_i=M_i\cdot \mathbf{v}.
\]

如果某个属性集合 $S$ 满足访问策略，则存在行集合 $I$ 和一组系数 $\omega_i$，使得：

\[
    \sum_{i\in I}\omega_i M_i=(1,0,\dots,0).
\]

于是秘密可以被重构：

\[
    \sum_{i\in I}\omega_i\lambda_i
    =
    \sum_{i\in I}\omega_i(M_i\cdot \mathbf{v})
    =
    (1,0,\dots,0)\cdot \mathbf{v}
    =
    s.
\]

这正是本项目解密判定的数学核心：能找到这样的 $\omega_i$ 就说明属性集合满足策略；找不到则说明属性集合不满足策略。

从实现角度看，解密时并不会重新遍历访问树来判断每个门是否满足，而是先根据用户属性筛选矩阵行，再求解一个线性方程组。这样做的好处是策略判断被统一为线性代数问题，重复属性、阈值门和普通 AND/OR 都可以走同一个求解接口。测试中的重复属性场景也正是为了验证这一点：多行可以映射到同一个属性，只要这些行被一起纳入方程组，就仍然可以求出重构系数。

Waters~[3]对 LSSS 的定义给出了本节公式的来源：矩阵 $M$ 是份额生成矩阵，$\rho(i)$ 指出第 $i$ 行属于哪个属性，向量 $(s,r_2,\dots,r_n)$ 生成所有行份额；当属性集合满足访问结构时，存在一组 $\omega_i$ 可以线性恢复秘密 $s$。实验三课件~[1]中的 Waters 2011 加密步骤同样先选取 $\mathbf{v}=(s,y_2,\dots,y_n)$，再计算每一行份额 $\lambda_i=\mathbf{v}\cdot M_i$。本项目的 \verb|Encrypt| 和 \verb|Decrypt| 函数就是围绕这一定义组织的。

\section{从访问树到 LSSS 的直觉}

访问树到 LSSS 的转换可以理解为“把树上的秘密共享过程展开成矩阵”。根节点持有秘密 $s$，每个内部阈值门使用 Shamir 风格的多项式把父节点份额分发给子节点。因为多项式求值对系数是线性的，所以每个叶子的最终份额都可以表示为 $s$ 和若干随机变量的线性组合。收集这些线性组合的系数，就得到 LSSS 矩阵。

Liu、Cao 和 Wong~[4]关注的问题是：LSSS 矩阵不如布尔公式或访问树直观，但高表达力 CP-ABE 又需要 LSSS 作为实现形式。该文提出从阈值访问树直接生成 LSSS 矩阵的思路，并指出对阈值门直接构造矩阵比先展开为 AND/OR 布尔公式更紧凑。本文没有逐字复现其算法，而是采用同样的思想：在树上递归传播线性份额向量，遇到阈值门时引入 $k-1$ 个随机维度，最终把叶子份额向量收集为矩阵行。

对示例策略 $(A\wedge B)\vee C$，OR 门不引入额外随机变量，AND 门引入一个随机变量。最终矩阵为：

\[
M=
\begin{bmatrix}
1 & 1\\
1 & 2\\
1 & 0
\end{bmatrix},
\quad
\rho=(A,B,C).
\]

其中 A、B 两行可以组合出 $(1,0)$：

\[
2(1,1)-1(1,2)=(1,0),
\]

因此 $\{A,B\}$ 可以恢复秘密。C 对应的行本身就是 $(1,0)$，所以 $\{C\}$ 也可以恢复秘密。单独的 A 行为 $(1,1)$，无法表示 $(1,0)$，因此 $\{A\}$ 不能恢复秘密。

这个例子也是后续实验场景设计的来源。场景一验证 $\{A,B\}$ 能通过两行线性组合恢复秘密，场景二验证 $\{C\}$ 能通过单行恢复秘密，场景三验证 $\{A\}$ 无法恢复秘密。也就是说，实验并不是随意选择三个属性集合，而是直接对应矩阵重构的三种典型情况：多行满足、单行满足和不满足。

\section{双线性对、指数抽象模型与真实配对实现}

CP-ABE 通常基于双线性对构造。设 $G_1$、$G_2$ 和 $G_T$ 是阶为素数 $p$ 的群，双线性对为：

\[
    e:G_1\times G_2\rightarrow G_T.
\]

其关键性质是双线性：

\[
    e(g_1^a,g_2^b)=e(g_1,g_2)^{ab}.
\]

真实 CP-ABE 方案的安全性依赖双线性群中的困难问题，例如 BDH、DBDH 或 Waters 2011 方案~[3]中更具体的安全假设。为了在课程实验中先把访问结构和代数抵消关系讲清楚，默认实现使用指数抽象模型：群元素 $g^x$ 直接存储为指数 $x$，群乘法对应指数加法，配对运算对应指数乘法。该模型使公式、矩阵和代码之间的关系更直接，适合验证 LSSS 访问结构和 CP-ABE 四算法流程。

在此基础上，当前 \verb|Project| 中的真实 pairing 路线进一步接入真实 pairing 后端。该后端使用 Python 的 \verb|py-ecc| 库，曲线选择 BLS12-381，标量域为曲线阶 $\mathbb{Z}_p$，群元素分布在 $G_1$、$G_2$ 和 $G_T$ 中。属性映射不再是简单 SHA-256 后取模，而是调用 \verb|hash_to_G1| 将属性字符串映射到 $G_1$；配对计算由 \verb|pairing| 函数完成；密钥和密文 JSON 中保存的也不再是单个整数指数，而是群元素坐标或 $G_T$ 的 12 个系数。

\begin{table}[H]
\centering
\caption{真实双线性群、指数抽象模型与 pairing 实现的对应关系}
\begin{tabular}{>{\raggedright\arraybackslash}p{0.28\textwidth} >{\raggedright\arraybackslash}p{0.28\textwidth} >{\raggedright\arraybackslash}p{0.30\textwidth}}
\toprule
真实双线性群表达 & 指数抽象表示 & BLS12-381 实现 \\
\midrule
$g^x\cdot g^y=g^{x+y}$ & 指数相加：$x+y \bmod p$ & 调用群加法/标量乘法生成真实 $G_1$ 或 $G_2$ 点 \\
$(g^x)^a=g^{ax}$ & 指数相乘：$a x \bmod p$ & 调用 \verb|multiply(point,a)| 完成标量乘法 \\
$e(g_1^x,g_2^y)=e(g_1,g_2)^{xy}$ & 配对结果指数为 $xy \bmod p$ & 调用 \verb|pairing(G2,G1)| 产生 $G_T$ 元素 \\
$e(g_1,g_2)^x\cdot e(g_1,g_2)^y$ & 目标群指数相加：$x+y \bmod p$ & 在 $G_T$ 中执行乘法和逆元运算 \\
\bottomrule
\end{tabular}
\end{table}

指数抽象模型能验证算法代数是否正确，但不能提供真实密码学安全性。真实 pairing 后端则把同一套访问树、LSSS 矩阵和重构逻辑放到 BLS12-381 群元素上运行，能够检查实现边界是否足够清晰：策略层和线性代数层保持稳定，变化主要集中在群元素生成、配对计算、哈希到群、序列化和密钥派生上。

需要强调的是，接入真实群运算并不自动等同于完成论文级安全证明。真实 CP-ABE 安全性还要求算法分量、随机数采样、哈希到群、序列化检查和安全模型声明与选定论文严格一致。本文的真实 pairing 路线主要用于说明底层运算已经从“整数指数模拟”推进到“真实曲线群元素计算”，并通过实验验证访问控制语义在真实配对路径下仍然成立。

\section{对称加密与完整性保护}

CP-ABE 通常不直接加密大文件，而是保护一个随机会话密钥，再由该会话密钥加密实际数据。这种设计称为混合加密。本项目也采用类似思路：指数抽象版随机选取一个消息指数 $m$，用其派生对称密钥；真实 pairing 版则随机生成 $G_T$ 中的会话元素，并由该 $G_T$ 元素的规范编码派生对称密钥。随后系统使用 AES-CBC 加密明文，并使用 HMAC-SHA256 保护对称加密载荷的完整性。

采用混合加密有两个原因。第一，ABE 运算适合保护较短的密钥材料，不适合直接处理任意长度文件；第二，对称加密在文件载荷处理上效率更高，也更容易与完整性校验组合。本项目中的 CP-ABE 部分负责保护会话密钥材料，对称加密部分负责保护真实明文文件。只有当用户属性集合满足访问策略并恢复出正确的消息指数或 $G_T$ 会话元素时，才能派生出正确的对称密钥并解开载荷。

密文 JSON 中的对称载荷保存在 \verb|payload| 字段下，主要包括三个子字段：\verb|nonce|、\verb|data| 和 \verb|tag|。其中 \verb|nonce| 在本实现中保存的是 AES-CBC 使用的 16 字节 IV（Initialization Vector，初始化向量）；\verb|data| 保存 AES-CBC 加密后的密文字节；\verb|tag| 保存 HMAC-SHA256 认证标签。这里使用字段名 \verb|nonce| 是为了表示“每次加密都重新生成的一次性随机值”，但在 AES-CBC 语境下，它的具体作用就是 IV。

IV 本身不是密钥，也不需要保密。它的作用是参与 CBC 模式第一轮分组运算，使同一明文在不同加密过程中不会因为使用相同对称密钥而产生完全相同的密文开头。解密时必须使用同一个 IV 才能正确还原明文，因此 IV 需要和密文一起保存。CP-ABE 密文分量和对称载荷共同组成最终密文：前者决定谁能恢复密钥材料，后者决定文件内容是否能够被正确解密。

采用“先加密、再认证”的模式有两个好处。第一，解密前可以先验证 HMAC，避免对被篡改密文执行 AES 解密。第二，HMAC 的输入是 \verb|IV || ciphertext|，也就是把 IV 和 AES-CBC 密文一起认证；因此无论攻击者修改 \verb|payload.nonce|、\verb|payload.data| 还是 \verb|payload.tag|，解密端都能在认证标签检查阶段发现异常。指数抽象版和真实 pairing 版都保留了这一完整性保护边界：访问策略满足只能说明用户有资格恢复密钥材料；如果载荷认证标签错误，系统仍应拒绝输出明文。

需要区分的是，HMAC 完整性保护并不改变 CP-ABE 访问结构本身。访问控制由 LSSS 重构方程决定，对称载荷完整性由 HMAC 标签决定，二者处在不同层次。这样的分层设计使实验结果更容易解释：属性不足导致的失败对应访问策略不满足，认证标签错误导致的失败对应密文载荷被篡改。第五章中的密文篡改测试正是用于验证这一边界。

在实验结果中，密文篡改测试专门修改 \verb|payload.tag|，用于验证系统不会把错误认证标签的密文继续交给 AES 解密。这一场景虽然不属于 LSSS 矩阵构造本身，但它说明系统在处理实际文件时没有忽略密文完整性问题。

\section{本章小结}

本章围绕 LSSS-CP-ABE 实验所需的核心概念展开，形成了从访问控制语义到程序实现对象的理论链路。ABE 将解密权限从单一身份扩展为属性集合，CP-ABE 进一步把访问策略放入密文，使数据拥有者能够在加密阶段指定访问条件。访问控制树提供了直观的策略书写形式，LSSS 则把树形策略转换为矩阵和行属性映射，使“属性集合是否满足策略”可以被表达为线性重构问题。

在访问结构层面，本章强调了阈值门的统一作用。AND、OR 和一般 $k$-of-$n$ 门都可以被看作阈值访问结构，因而不必为不同布尔门分别设计矩阵构造规则。示例策略 $(A\wedge B)\vee C$ 展示了多行重构、单行重构和无法重构三种典型情况，并在第五章分别对应 $\{A,B\}$ 成功、$\{C\}$ 成功和 $\{A\}$ 失败。重复属性和 $\rho$ 非单射的讨论也说明，矩阵行与属性名之间并不总是一一对应，解密端必须根据行集合而不是简单属性去重来求重构系数。

在密码运算层面，真实 CP-ABE 方案依赖双线性群和困难问题假设。本实验先采用指数抽象模型表示群元素和配对结果，以便直接观察加解密公式中的指数抵消关系；随后在同一项目中接入 BLS12-381 pairing 后端，把 $G_1/G_2/G_T$ 群元素、哈希到群、配对和序列化纳入运行流程。前者有利于解释算法结构，后者有利于验证底层群运算真实化后的可运行性。对称加密与完整性保护则承担实际文件载荷处理：CP-ABE 层决定能否恢复密钥材料，AES-CBC-HMAC 层决定明文载荷能否被正确解密并通过篡改检测。

由此可见，本章的理论梳理也构成了第三章至第五章的实现约束。策略解析必须保留阈值门层次和重复叶节点，矩阵构造必须保持行向量与属性映射 $\rho$ 的对应关系，解密判定必须以线性重构系数是否存在为依据，密文处理则需要区分访问控制失败与载荷完整性失败。实验结果的有效性不能只由“是否输出明文”判断，还需要同时考察策略语义、矩阵重构、密钥材料恢复和对称载荷认证四个环节是否一致。

从参考资料对应关系看，Bethencourt 等人~[2]说明了 CP-ABE 访问控制模型，Waters~[3]说明了 LSSS 矩阵和重构系数如何进入 CP-ABE 算法，Liu 等人~[4]说明了阈值访问树到 LSSS 的工程转换价值，实验三课件~[1]则把这些概念整理为课程实现任务。因此，第三章的方案设计、第四章的系统实现和第五章的实验验证都以本章的概念链条为基础：第三章给出树到矩阵和四算法的数据流，第四章说明这些公式如何落到 PowerShell 代码，第五章验证理论语义与运行结果的一致性。
